\chapter*{Resumen}
\addcontentsline{toc}{chapter}{Resumen}

Las redes industriales OT que soportan procesos productivos y servicios esenciales se han convertido en un objetivo prioritario de la ciberseguridad moderna. Su protección no puede abordarse del mismo modo que una red IT convencional: la disponibilidad de la planta, la seguridad física, los equipos heredados, las ventanas de mantenimiento reducidas y la convivencia de protocolos industriales con tecnologías IP hacen que cualquier auditoría intrusiva, prueba de penetración o cambio de configuración sobre producción pueda generar impactos operativos relevantes.

Este Trabajo Fin de Máster se plantea en colaboración con AMC Global, grupo empresarial español del sector de alimentación y bebidas naturales con infraestructuras IT y OT distribuidas en varios centros productivos. La colaboración aporta un contexto industrial realista: una organización con procesos de producción continuos, necesidades de inventario técnico, exigencias de ciber-resiliencia y obligaciones crecientes vinculadas al Esquema Nacional de Seguridad, la Directiva NIS2 y marcos de referencia como ISA/IEC 62443 y NIST SP 800-82.

La propuesta de valor del trabajo consiste en explorar la viabilidad de aplicar gemelos digitales a la ciberseguridad de redes OT. Para ello se diseña la arquitectura de GEMEROTIC, un constructor de gemelos digitales orientado a redes industriales. La idea no es crear únicamente un diagrama de red, sino establecer las bases de una herramienta de modelado técnico, comparable conceptualmente a un entorno de diseño asistido, capaz de representar activos, ubicaciones, conexiones, segmentación, configuración y evidencias de validación.

La arquitectura propuesta separa la información en tres capas: infraestructura física, conectividad lógica y seguridad OT. Sobre este modelo se construye un flujo que parte de una interfaz visual, persiste el estado editable del proyecto, genera inventario enriquecido, produce artefactos reproducibles y permite desplegar o validar el gemelo en un entorno controlado. NetBox se emplea como repositorio derivado de inventario técnico, mientras que PostgreSQL mantiene el estado granular del gemelo dentro de GEMEROTIC. El pipeline genera artefactos para herramientas como Containerlab, Ansible, Batfish y OPA, incorporando además una capa inicial de cumplimiento determinista y asistencia explicativa basada en IA.

El resultado esperado es una base sólida para un constructor de gemelos digitales aplicable a redes OT, con capacidad para reproducir topologías existentes, analizar cambios antes de aplicarlos en planta, apoyar auditorías técnicas y facilitar procesos de mejora continua sin interrumpir la operación real. El trabajo no pretende certificar automáticamente el cumplimiento normativo ni sustituir una auditoría formal, sino demostrar una arquitectura viable para estructurar evidencias, reducir riesgo operativo y evolucionar hacia validaciones más avanzadas.

\chapter*{Abstract}
\addcontentsline{toc}{chapter}{Abstract}

Industrial OT networks supporting production processes and essential services have become a major cybersecurity concern. Their protection cannot be approached in the same way as a conventional IT network: plant availability, physical safety, legacy equipment, limited maintenance windows and the coexistence of industrial protocols with IP technologies make intrusive audits, penetration testing or configuration changes on production systems potentially disruptive.

This Master's Thesis is developed in collaboration with AMC Global, a Spanish business group in the natural food and beverage sector with distributed IT and OT infrastructures across several production sites. This collaboration provides a realistic industrial context: an organisation with continuous production processes, technical inventory needs, cyber-resilience requirements and increasing obligations related to the Spanish National Security Framework, the NIS2 Directive and reference frameworks such as ISA/IEC 62443 and NIST SP 800-82.

The value proposition of this work is to explore the feasibility of applying digital twins to OT network cybersecurity. To this end, the thesis designs the architecture of GEMEROTIC, a digital twin constructor focused on industrial networks. The objective is not merely to create a network diagram, but to establish the foundations of a technical modelling tool, conceptually similar to a computer-aided design environment, able to represent assets, locations, connections, segmentation, configuration and validation evidence.

The proposed architecture separates information into three layers: physical infrastructure, logical connectivity and OT security. Based on this model, the system starts from a visual interface, persists the editable project state, generates enriched inventory, produces reproducible artefacts and enables the digital twin to be deployed or validated in a controlled environment. NetBox is used as a derived technical inventory repository, while PostgreSQL stores the granular state of the twin within GEMEROTIC. The pipeline generates artefacts for tools such as Containerlab, Ansible, Batfish and OPA, also incorporating an initial deterministic compliance layer and AI-assisted explanation.

The expected result is a solid foundation for a digital twin constructor applicable to OT networks, capable of reproducing existing topologies, analysing changes before applying them to plant environments, supporting technical audits and enabling continuous improvement without interrupting real operations. The work does not aim to automatically certify regulatory compliance or replace a formal audit, but to demonstrate a viable architecture for structuring evidence, reducing operational risk and evolving towards more advanced validations.
